\subsection{Notation}
$L_2^n[a,b]$ denotes the set of $\R^n$-valued, Lebesgue square-integrable equivalence classes of functions on the spatial domain $[a,b]\subset\R$. The space $\R L_2^{\mathbf{n}}[a,b]$ denotes the Cartesian product $\R^{n_1} \times L_2^{n_2}[a,b]$, where $\mathbf{n} = \bmat{n_1 \\ n_2}$, with inner product
\begin{equation*}
	\ip{\bmat{x_1 \\ \mathbf{x}_2}}{\bmat{y_1 \\ \mathbf{y}_2}}_{\R L_2}
	:= x_1^\top y_1 + \ip{\mathbf{x}_2}{\mathbf{y}_2}_{L_2}.
\end{equation*} When clear from context, we occasionally omit the subscript and simply write $\ip{\cdot}{\cdot} := \ip{\cdot}{\cdot}_{\R L_2}$. Since the only spatial domain considered here is $[a,b]$, we typically omit the domain and write $L_2^n$ or $\R L_2^{\mathbf{n}}$, further omitting dimensions when clear from context. For a normed space $X$ and Banach space $Y$, $\mcl L(X,Y)$ denotes the Banach space of bounded linear operators from $X$ to $Y$ with induced norm
\[
	\|\mcl P\|_{\mcl L(X,Y)} = \sup_{\|\mbf{x}\| = 1} \|\mcl P \mathbf{x}\|_Y,
\] and $\mcl L(X) := \mcl L(X,X)$. Our notational convention is to write functions in \textbf{bold} (e.g., $\mathbf{x}$, indicating $\mathbf{x} \in \RL$) and operators in calligraphic font (e.g., $\mcl P \in \mcl L(\RL)$).
For a Hilbert space $X$ and operator $\mcl A \in \mcl L(X)$, the adjoint $\mcl A^*$ is defined by
\[
	\ip{\mathbf{x}}{\mcl A \mathbf{y}}_X = \ip{\mcl A^* \mathbf{x}}{\mathbf{y}}_X,
	\quad \forall\, \mathbf{x},\mathbf{y} \in X.
\] The space $L_2^\mbf{n}[0,\infty)$ denotes $\R^n$-valued square-integrable signals defined on the temporal domain $[0,\infty)$,
 so that $x(t) \in \R^n$. Similarly, $\RL^\mbf{n}[0,\infty)$ denotes $\RL$-valued square-integrable temporal signals. 
 For suitably differentiable $\mathbf{x} \in \RL^\mbf{n}[0,\infty)$, we write $\mathbf{\dot{x}} := \frac{\partial \mathbf{x}}{\partial t}$. For a Hilbert space $\mathbf{H}$, we denote by $\mathbf{L}_{\mathbf{H}}$ the space of square-integrable functions $u : [0,\infty) \to \mathbf{H}$ with inner product
\[
	\ip{\mathbf{u}}{\mathbf{v}}_{\mathbf{L}_{\mathbf{H}}}
	= \int_0^\infty \ip{\mathbf{u}(t)}{\mathbf{v}(t)}_{\mathbf{H}}\, dt.
\] We define the extended space $\mathbf{L}_{e,\mathbf{H}}$ to consist of functions $u : [0,\infty) \to \mathbf{H}$ such that
\[
	\int_0^T \|\mathbf{u}(t)\|_{\mathbf{H}}^2\, dt < \infty
	\quad \forall\, T \ge 0.
\] We use $\mathbf{L}_{e,[a,b]}$ and $\mathbf{L}_{[a,b]}$ to denote $\mathbf{L}_{e,\mathbf{H}}$ and $\mathbf{L}_{\mathbf{H}}$ with $\mathbf{H} = \R L_2[a,b]$. For any Hilbert space $\mathbf{H}$, the truncation operator $p_\tau : \mathbf{L}_{e,\mathbf{H}} \to \mathbf{L}_{e,\mathbf{H}}$ is defined for $\mathbf{y} \in \mathbf{L}_{e,\mathbf{H}}$ by
\begin{equation*}
	(p_\tau \mathbf{y})(t) =
	\begin{cases}
		\mathbf{y}(t), & 0 \le t \le \tau, \\
		0,             & \text{otherwise}.
	\end{cases}
\end{equation*}
\begin{definition}
	Let $\mathbf{H}$ and $\mathbf{G}$ be Hilbert spaces. For all $\tau\geq0$, an operator $G : \mathbf{L}_{\mathbf{H}} \to \mathbf{L}_{\mathbf{G}}$ is:
	\begin{itemize}
		\item \textbf{Causal} if $p_\tau G = p_\tau G p_\tau$.
		\item \textbf{Bounded on $\mathbf{L}_{\mathbf{H}}$} if $\|G \mathbf{v}\|_{\mathbf{L}_{\mathbf{G}}} \le C \|\mathbf{v}\|_{\mathbf{L}_{\mathbf{H}}}$ for some $C \ge 0$ and all $\mathbf{v} \in \mathbf{L}_{\mathbf{H}}$.
		\item \textbf{Bounded on $\mathbf{L}_{e,\mathbf{H}}$} if $\|p_\tau G \mathbf{v}\|_{\mathbf{L}_{\mathbf{G}}} \le C \|p_\tau\mathbf{v}\|_{\mathbf{L}_{\mathbf{H}}}$ for some $C \ge 0$ and all $\mathbf{v} \in \mathbf{L}_{e,\mathbf{H}}$.
		      % \item \textbf{Incrementally $\mathbf{L}_{e,\mathbf{H}}$-bounded} if for some $C \geq 0$,
		      % \[
		      % \|p_\tau (G \mathbf{v} - G \mathbf{u})\|_{\mathbf{L}_{\mathbf{G}}}
		      % \le C \|p_\tau (\mathbf{v} - \mathbf{u})\|_{\mathbf{L}_{\mathbf{H}}}
		      % \]
		      % for all $\mathbf{v}, \mathbf{u} \in \mathbf{L}_{e,\mathbf{H}}$.
	\end{itemize}
\end{definition}
Finally, for a given $\mcl K \in \mcl L(\RL^\mbf{n}, \RL^\mbf{m})$, we define the operator $\mcl M_{\mcl K} : \mathbf{L}_{e,[a,b]}^\mbf{n} \to \mathbf{L}_{e,[a,b]}^\mbf{m}$ to compactly represent a multiplication operator. For $\mathbf{u} = \bmat{u_1 & \cdots & u_N}^\top$ with $\mathbf{u} \in \RL$, its action at time $t$ is given by the block structure of $\mcl K_{ij}$:
\begin{equation}
	\label{eq:block_operator_def}
	(\mcl M_{\mcl K} \mathbf{u})(t)
	:= \bmat{
		\mcl K_{11} & \cdots & \mcl K_{1n} \\
		\vdots      & \ddots & \vdots      \\
		\mcl K_{m1} & \cdots & \mcl K_{mn}
	}
	\bmat{\mbf{u}_1 \\ \vdots \\ \mbf{u}_n}(t).
\end{equation}

\subsection{PI operators}
Partial Integral operators, also known as PI operators, are a class of bounded linear integral operators that are defined jointly on a vector space and Hilbert space.
\begin{definition}
	We say $\mcl P \in {\Pi}_4^{\mbf{m}\times\mbf{n}} \subset \mcl L(\RL^\mbf{n}, \RL^\mbf{m})$, denoted by $\prod \begin{bmatrix}
			P   & Q_1   \\
			Q_2 & \{R\} \\
		\end{bmatrix}$, if there exist a matrix $P$ and matrix polynomials $Q_1, Q_2, R_0, R_1$, and $R_2$ (of compatible dimension) such that,
	\begin{gather*}
		\Bigg[\Bigg(\prod \begin{bmatrix}
				P   & Q_1   \\
				Q_2 & \{R\} \\
			\end{bmatrix}\Bigg)\! \! \bmat{x \\ \mbf{x}}\Bigg](s)\! :=\! \bmat{Px\! +\!\int\limits_a^b\!Q_1(s)\mbf x(s)ds\\Q_2(s)x\!+\!\mcl R \mbf x(s)},\\
		(\mcl R \mbf x)(s)\!\! =\!\! R_0(s)\mbf x(s)\! +\! \int\limits_a^s\!\!\!R_1(s,\theta)\mbf x(\theta)d\theta\! +\!\int\limits_s^b\!\!\!R_2(s,\theta)\mbf x(\theta) d\theta.
	\end{gather*}

\end{definition}
The vector space of PI operators, given compatible dimensions, are closed under composition, addition, adjoint and concatenation -- implying that ${\Pi}_4$ is a composition algebra \cite{shivakumar2024GDPE}.

\subsection{Linear PI Inequalities (LPIs)}
\begin{definition}(LPIs)
	\label{def_LPI}
	For given PI operators $\{\mcl E_{i,j}, \mcl{F}_{i,j}, \mcl{G}_i\}$ and a convex linear functional $\mcl L(\cdot)$, a linear PI inequality is a convex optimization of the following form
	\begin{align}
		 & \min\limits_{P_i, Q_{1i}, Q_{2i}, R_{0i}, R_{1i}, R_{2i}}  \mcl L(\{P_i, Q_{1i}, Q_{2i}, R_{0i}, R_{1i}, R_{2i}\})\notag \\
		 & \mathrm{s. t. }\ \sum\limits_{j = 1}^K \mcl E_{ij}^*\prod \begin{bmatrix}
			                                                             P_i    & Q_{1i}  \\
			                                                             Q_{2i} & \{R_i\} \\
		                                                             \end{bmatrix}\mcl F_{ij} +\mcl G_i \succeq 0
	\end{align}
\end{definition}
Only LMIs are required to solve an optimization problem involving LPIs.
\subsection{PIETOOLS}
PIETOOLS is a MATLAB toolbox used for declaration, manipulation and solving LPIs \cite{9147712}. For learning how to use PIETOOLS and its functionalities visit \href{https://control.asu.edu/pietools/pietools}{PIETOOLS homepage}.

